%==============================================================================
% Digital Twin rApp - Downlink Throughput Engine : Build Report
%
% HOW TO USE THIS TEMPLATE IN OVERLEAF
% 1. Create a new Overleaf project and paste this file as main.tex.
% 2. Upload your 8 images using the exact file names referenced below
%    (fig_geojson.png, fig_3d_scene.png, ...). PNG/JPG/PDF all work.
% 3. Flip the switch on the next line from \useplaceholderstrue to
%    \useplaceholdersfalse. Every placeholder box then becomes your real image.
%==============================================================================
\documentclass[11pt]{article}

\usepackage[T1]{fontenc}
\usepackage[utf8]{inputenc}
\usepackage[margin=1in]{geometry}
\usepackage{graphicx}
\usepackage{float}
\usepackage{enumitem}
\usepackage{underscore} % lets us type underscores in text without escaping
\usepackage{xcolor}
\usepackage{hyperref}

% ---- image placeholder switch -----------------------------------------------
\newif\ifuseplaceholders
\useplaceholderstrue % <-- change to \useplaceholdersfalse once images are uploaded

% \figobs{filename}{width-fraction}{caption}{observation}
\newcommand{\figobs}[4]{%
  \begin{figure}[H]
    \centering
    \ifuseplaceholders
      \fbox{\parbox[c][4.5cm][c]{#2\textwidth}{\centering\itshape [\,Image placeholder:\ \texttt{#1}\,]}}%
    \else
      \includegraphics[width=#2\textwidth]{#1}%
    \fi
    \caption{#3}
  \end{figure}%
  \noindent\textbf{Observation.}~#4\par\medskip
}

% ---- title ------------------------------------------------------------------
\title{Digital Twin rApp\\[2pt]\large Downlink Throughput Engine --- Build Report}
\author{[Your Name] \\ [Team / Group]}
\date{[Date]}

\begin{document}
\maketitle

\begin{abstract}
I have built the downlink throughput engine for the Digital Twin rApp. It takes a
real-world area, reconstructs it as a ray-traced 3D radio environment, and
computes downlink throughput per UE and per cell.
\end{abstract}

%==============================================================================
\section{Two Distinct Parts --- Please Don't Conflate Them}

\subsection*{A) The Engine (\texttt{dtrapp} package + CLI)}
This is the actual product. It runs the whole simulation and produces the results
in the output directory, from one command:
\begin{verbatim}
python3 -m dtrapp.runner.cli configs/example.yaml
\end{verbatim}
It needs nothing else to run.

\smallskip
\noindent\textbf{Note:} It will not run on older GPUs present on the remote Linux
provided to us. Please set the following before running the CLI. It runs
reasonably fast on CPU for this small scene.
\begin{verbatim}
export CUDA_VISIBLE_DEVICES=""
\end{verbatim}

\subsection*{B) The Notebook (\texttt{dtrapp\_sionna\_walkthrough.ipynb})}
Purely for visualization and understanding. It re-runs the same pipeline step by
step with plots so anyone can see what's happening. It is not part of the run
path --- the engine runs fine without it.

%==============================================================================
\section{Abstract --- Pipeline End to End (The Engine)}

The only input is one \texttt{config.yaml}. Its primary constituent is a lat/lon
bounding box, given as its two diagonal corners (\texttt{min_lat}/\texttt{min_lon},
\texttt{max_lat}/\texttt{max_lon}). The same file sets:

\begin{enumerate}[leftmargin=*]
  \item \textbf{Geometry:} \texttt{default_building_height_m} (fallback when OSM
  lacks a height), \texttt{meters_per_level}, \texttt{building_material}.
  \item \textbf{Network} (seeded RNG stand-in for live data): \texttt{seed} (full
  reproducibility), \texttt{num_sites}, \texttt{sectors_per_site},
  \texttt{bs_height_m}, \texttt{tx_power_dbm}, \texttt{carrier_freq_hz},
  \texttt{bandwidth_hz}, \texttt{num_ues}, \texttt{ue_height_m},
  \texttt{ue_noise_figure_db}.
  \item \textbf{Antennas} (configurable, not hardcoded):
  \texttt{bs_antenna_rows/cols/pattern/polarization} (e.g., a $4\times1$ TR-38.901
  vertically-polarized panel), the \texttt{ue_antenna_*} equivalents (e.g.,
  $1\times1$ isotropic), \texttt{antenna_spacing} (in wavelengths), and BS
  \texttt{downtilt_deg}.
  \item \textbf{KPI / Link Layer:} \texttt{temperature_k} (thermal-noise
  reference), \texttt{bler_target} (link-adaptation BLER target),
  \texttt{mcs_table_index} (5G-NR MCS table 1/2), and the OFDM grid
  (\texttt{subcarrier_spacing_hz}, \texttt{num_subcarriers},
  \texttt{num_ofdm_symbols}).
  \item \textbf{Output Directory.}
\end{enumerate}

From the bbox, the engine queries the OpenStreetMap Overpass API for building
footprints, extrudes them into a 3D mesh, and writes a Mitsuba scene that Sionna
RT loads --- i.e., it reconstructs the real environment in the ray tracer (no
fallback geometry: if OSM returns nothing, it errors out rather than fabricating a
scene). It places one transmitter per cell and one receiver per UE, ray-traces the
channel frequency response (CFR), and from that computes per-UE SINR and then
downlink throughput (Mbps), written to \texttt{ue_throughput.csv},
\texttt{cell_throughput.csv}, and \texttt{throughput.json}.

%==============================================================================
\section{Swappable Data Layer for Real Network Data}

The cell/UE generator sits behind a \texttt{NetworkDataSource} interface; the rest
of the pipeline depends only on that interface. Integrating live network data means
implementing one class (projecting real BS/UE coordinates into the scene's local
ENU frame) --- no downstream changes.

%==============================================================================
\section{Sionna SYS Link Model (The Throughput Stage)}

\begin{itemize}[leftmargin=*]
  \item Sionna RT yields the CFR.
  \item Sionna SYS turns it into a realistic 5G-NR rate:
  \begin{itemize}
    \item Post-equalization SINR via regularized zero-forcing (RZF) precoding ---
    capturing the BS array's beamforming gain --- plus an LMMSE equalizer.
    Inter-cell interference from neighbouring cells is included (folded into the
    effective noise), not just thermal noise.
    \item Link adaptation (\texttt{InnerLoopLinkAdaptation} +
    \texttt{PHYAbstraction}) selects the highest MCS whose BLER stays within
    \texttt{bler_target}, using NVIDIA's NR BLER tables $\rightarrow$ a
    modulation-capped spectral efficiency (throughput is bounded by the real MCS
    ceiling, not unbounded Shannon).
    \item Proportional-fair scheduling shares each cell's airtime across its
    attached UEs.
  \end{itemize}
\end{itemize}

%==============================================================================
\section{Visualization \& Verification (via the Notebook \& geojson.io)}

\subsection{GeoJSON / Overpass Turbo Cross-check}
The bbox polygon and building footprints match the real map, and each scene shape
maps back to its OSM way ID, confirming the reconstruction is faithful.

\figobs{fig_geojson.png}{0.8}
{geojson.io / Overpass Turbo: the bbox and building footprints over the real map.}
{The bounding-box polygon and the extracted building footprints align with the real
map, and each footprint maps back to its OSM way ID. This confirms the
reconstruction is a faithful copy of the OpenStreetMap data for the chosen area.}

\subsection{Per-stage Visualization (The Notebook)}
The notebook is where the results are made visible and checkable, stage by stage:
the 3D scene, the BS/UE layout, the ray paths, the path-gain heatmap, and the
SINR/throughput distributions.

\figobs{fig_3d_scene.png}{0.8}
{The 3D scene reconstructed in Sionna RT.}
{The OSM footprints are extruded into flat-roof prisms over a flat ground plane and
loaded into Sionna RT as a Mitsuba scene. This is the ray-traced 3D reconstruction
of the real environment in which propagation is computed.}

\figobs{fig_bs_ue_positions.png}{0.8}
{Base-station and UE positions with their IDs.}
{Top-down layout of the seeded network, with one transmitter per cell and one
receiver per UE, each labelled by ID. Sites are placed uniformly with sectors split
evenly ($360^\circ/N$), and UEs are placed uniformly across the scene.}

\figobs{fig_rt_paths.png}{0.8}
{Ray-traced propagation paths between transmitters and receivers.}
{The actual rays traced between transmitters and receivers --- the LOS path plus up
to \texttt{max_depth}~$=3$ reflections (no diffraction). It shows how the building
geometry channels and blocks propagation.}

\figobs{fig_path_gain.png}{0.8}
{Per-link path gain for every BS--UE pair.}
{Ray-traced path gain for each BS--UE link. Strong gains correspond to LOS / nearby
links, while weak gains correspond to NLOS links shadowed by buildings --- consistent
with the scene geometry.}

\figobs{fig_sinr_throughput_dist.png}{0.85}
{Distribution of per-UE SINR and downlink throughput.}
{Distribution of post-equalization SINR and the resulting downlink throughput across
all UEs. Throughput is capped by the selected 5G-NR MCS (modulation-limited, not
unbounded Shannon).}

\subsection{Physically-interpretable Maps}
Top-down plots overlay buildings, cells (with sector boresight) and UEs coloured by
throughput, so results are explainable: a UE in NLOS behind a building shows low
SINR and low throughput, as expected.

\figobs{fig_ue_throughput_map.png}{0.85}
{UEs coloured by throughput, each linked to its serving cell.}
{Top-down map overlaying buildings, cells (with sector boresight) and UEs coloured by
throughput, each linked to its serving cell. A UE in NLOS behind a building shows low
SINR and low throughput, exactly as expected.}

\subsection{Two-UE Factor Comparison}
For any UE pair, a side-by-side table (distance, path gain, interference, SINR, cell
load, bandwidth share, throughput) decomposes the throughput gap into resource
sharing (cell load) vs.\ SINR and names the dominant driver.

\figobs{fig_ue6_ue7_comparison.png}{0.9}
{Qualitative comparison of UE6 vs.\ UE7 (engine output).}
{Side-by-side decomposition for the UE pair separating the throughput gap into
resource-sharing vs.\ SINR effects. Here a closer LOS UE underperforms because it sits
near a cell edge with high inter-cell interference.}

%==============================================================================
\section{Status --- Remaining}
\begin{enumerate}[leftmargin=*]
  \item Ingest real network data via the swappable layer.
  \item Calibrate against a reference so absolute throughput matches a real network.
  \item Per-device heterogeneous antennas (Sionna RT applies one TX/RX array per
  solve, so this needs a grouped-per-antenna-type solve).
\end{enumerate}

%==============================================================================
\section{Open Problem --- Validating Absolute Throughput}
I can verify the engine's behaviour (modulation-capped rates, per-cell resource
sharing, SINR/throughput consistent with geometry and interference), but I cannot
yet validate the absolute Mbps values, because the network is randomly generated and
there's no ground-truth measurement to compare against. Closing this needs either
real network data or a trusted calibration reference.

\medskip
\noindent\textbf{Code availability.} The entire codebase is uploaded to datastore:
\begin{verbatim}
\\datastore\network_team\Network_Team\General\interns\btech_26\
Digital_Twin_krushnansh\minsysdtrapp
\end{verbatim}
Please run the \texttt{.ipynb} file to get an easy understanding of the codebase.

%==============================================================================
\section{Assumptions}

\subsection{Environment / Geometry}
\begin{itemize}[leftmargin=*]
  \item Buildings only, flat ground. Scene = OSM building footprints extruded to
  flat-roof prisms over a flat ground plane at $z=0$. No terrain/elevation, no
  vegetation, vehicles, street furniture, or indoor structure.
  \item Single radio material for everything. All buildings and the ground share one
  material (e.g., \texttt{itu_concrete}); no glass/brick/metal differentiation.
  \item Building heights are often inferred: explicit OSM height, else number of
  levels $\times$ \texttt{meters_per_level}, else a default value (many buildings are
  guessed).
  \item Local flat-earth (equirectangular) projection about the bbox centre.
\end{itemize}

\subsection{Antennas}
\begin{itemize}[leftmargin=*]
  \item Uniform arrays: one BS antenna array applies to ALL base stations and one UE
  array to ALL UEs (Sionna RT uses a single TX/RX array per solve). No per-device /
  heterogeneous antennas.
  \item Fixed electrical downtilt for every sector; fixed element spacing
  (0.5 wavelength by default).
\end{itemize}

\subsection{Network Data (Current Stand-in)}
\begin{itemize}[leftmargin=*]
  \item Randomly generated, not real --- seeded RNG for BS/UE positions and config.
  \item Single carrier frequency for all cells $\rightarrow$ full frequency reuse
  (reuse-1): every other cell is a co-channel interferer. No multi-carrier / carrier
  aggregation.
  \item Homogeneous cells: identical Tx power, bandwidth, and antenna for every cell.
  \item Random uniform placement of sites (sectors split evenly, $360^\circ/N$) and
  UEs; UEs are outdoor at a fixed height and may even fall inside a building footprint
  (no placement check).
  \item Static, single snapshot: no UE mobility, no time evolution. Generated traffic
  demand is unused (full-buffer assumption).
\end{itemize}

\subsection{Propagation (Sionna RT)}
\begin{itemize}[leftmargin=*]
  \item Bounded ray interactions: \texttt{max_depth} reflections (default 3); no
  diffraction modelling.
  \item Static channel, perfect CSI: no fast fading / Doppler, and channel knowledge
  is assumed perfect (no estimation error).
\end{itemize}

\subsection{SINR / Link Layer (Sionna SYS)}
\begin{itemize}[leftmargin=*]
  \item Inter-cell interference = full-buffer neighbours folded into the noise: every
  other cell is assumed to transmit at full power on all resources (worst-case,
  constant interference), treated as added noise rather than joint detection.
  \item Single-stream SU-MIMO: one UE per resource element per cell, so no intra-cell
  multi-user interference and no MIMO spatial-multiplexing layers (the BS array
  provides beamforming gain only).
  \item Representative sub-band SINR: post-equalization SINR is computed on a small
  OFDM grid (\texttt{num_subcarriers}, e.g., 128) and assumed representative;
  throughput then scales with the full configured \texttt{bandwidth_hz}.
  \item Frequency-flat MCS selection: a single effective SINR per UE drives the MCS
  choice.
\end{itemize}

\subsection{Throughput / Scheduling}
\begin{itemize}[leftmargin=*]
  \item Proportional-fair = equal airtime: valid for this static, full-buffer single
  snapshot (PF reduces to $1/K$ per cell).
  \item BLER pinned at target: goodput = spectral efficiency $\times$
  \texttt{(1 - bler_target)}; no explicit control-channel / reference-signal overhead
  beyond the PHY abstraction.
  \item Strongest-cell association: each UE attaches to the single highest-received-
  power cell; no load balancing, CoMP, or handover.
  \item Downlink only.
  \item Single MCS table (64QAM by default; 256QAM optional).
\end{itemize}

\subsection{Validation}
\begin{itemize}[leftmargin=*]
  \item Uncalibrated: the model is not yet calibrated against real measurements, so
  absolute throughput values are unvalidated.
\end{itemize}

\end{document}
